%%  MIE 1626, Summer 2026 -- Test reference sheet
%%  The test is open-handout (this sheet is distributed in advance), not open-book.
%%  Listed here is the R syntax used in the course, so students don't have to memorize it.

\documentclass[10pt]{article}
\usepackage[hmargin=.5in,vmargin=.6in]{geometry}
\usepackage{multicol}
\usepackage{fancyhdr}
\usepackage{lastpage}
\usepackage{titlesec}

% Compact section formatting
\titleformat{\section}{\normalsize\bfseries\scshape}{}{0pt}{}
\titlespacing{\section}{0pt}{1.5ex}{0.5ex}

% Headers / footers
\pagestyle{fancy}
\renewcommand{\headrulewidth}{0pt}
\renewcommand{\footrulewidth}{0pt}
\lhead{MIE 1626}
\chead{\scshape Test Reference Sheet}
\rhead{Summer 2026}
\lfoot{\small Department of Mechanical and Industrial Engineering, University of Toronto}
\cfoot{}
\rfoot{\small Page \thepage\ of \pageref{LastPage}}

\setlength{\parindent}{0pt}
\setlength{\parskip}{0.25ex}
\setlength{\columnsep}{1.5em}

% Each entry: a function signature in tt, indented description on the next line.
\newcommand{\entry}[2]{%
  \par\medskip\noindent\textbf{\texttt{#1}}\par%
  \hspace*{1.5em}\parbox[t]{\dimexpr\linewidth-1.5em}{#2}\par%
}

\begin{document}

\begin{center}
  {\large\bfseries MIE 1626 --- Summer 2026 --- Test Reference Sheet}
\end{center}
\smallskip

This sheet is distributed in advance and is the only material allowed at the test
(the test is \emph{open handout}, not open book). It lists the R syntax used in
the course so that you do not need to memorize it. You should still know
\emph{what} each function does and \emph{when} to use it; the sheet only spares
you having to remember the spelling.

\begin{multicols}{2}

\section*{Functions and control flow}
\entry{f <- function(x, y) \{ ... return(...) \}}{Define a function. The body is in curly braces; the returned value comes from \texttt{return()} or, if there is no explicit \texttt{return}, the value of the last expression.}
\entry{if (cond) \{ ... \} else \{ ... \}}{Conditional.}

\section*{Vectors}
\entry{c(1, 2, 3)}{Build a vector.}
\entry{v[3]}{One element (R is 1-indexed).}
\entry{v[2:5]}{A range.}
\entry{v[v > 0]}{Boolean mask: keep the elements where the mask is TRUE.}
\entry{\&, |, !}{Element-wise and, or, not.}
\entry{sapply(v, f)}{Apply \texttt{f} to every element of \texttt{v}; return a vector.}
\entry{replicate(n, expr)}{Run \texttt{expr} \texttt{n} times; return the results.}
\entry{mean(v), sum(v), length(v)}{Standard reductions.}
\entry{which.max(v)}{Index of the largest element.}

\section*{Random sampling and simulation}
\entry{sample(x, size, replace = FALSE)}{A random sample of \texttt{size} values drawn from \texttt{x}, with or without replacement.}
\entry{rbinom(n, size, prob)}{\texttt{n} draws from a Binomial distribution with \texttt{size} trials and success probability \texttt{prob}. For a single Bernoulli trial, use \texttt{size = 1}.}
\entry{rnorm(n, mean = 0, sd = 1)}{\texttt{n} draws from a Normal distribution.}
\entry{plogis(z)}{The logistic / sigmoid function: \texttt{1 / (1 + exp(-z))}. Inverse of the logit.}

\section*{Data frames}
\entry{df\$col}{One column as a vector.}
\entry{df[row, col]}{Subset by row and column.}
\entry{df[df\$x > 0, ]}{Subset the rows by a boolean mask. Note the trailing comma.}
\entry{nrow(df), ncol(df)}{Row and column counts.}
\entry{head(df), tail(df)}{First or last rows.}

\section*{The pipe and the tidyverse}
\entry{x \%>\% f()}{The pipe: equivalent to \texttt{f(x)}. Chains compose: \texttt{x \%>\% f() \%>\% g()} is \texttt{g(f(x))}.}
\entry{filter(df, cond)}{Keep rows where \texttt{cond} is TRUE.}
\entry{select(df, col1, col2)}{Keep only the named columns.}
\entry{arrange(df, col)}{Sort rows ascending. \texttt{arrange(df, desc(col))} sorts descending.}
\entry{rename(df, new = old)}{Rename a column.}
\entry{mutate(df, new\_col = expr)}{Add or replace a column.}
\entry{group\_by(df, col)}{Group rows for the next \texttt{summarize} or \texttt{mutate}.}
\entry{summarize(df, stat = f(col))}{Collapse a (possibly grouped) data frame to one row per group.}
\entry{distinct(df, col)}{Unique rows by column.}
\entry{n\_distinct(col)}{Count of unique values.}

\section*{Linear regression}
\entry{lm(y \textasciitilde{} x1 + x2, data = df)}{Fit a linear regression of \texttt{y} on \texttt{x1} and \texttt{x2}.}
\entry{summary(model)}{Print the regression summary: coefficient estimates, standard errors, t-values, p-values, residual standard error, R-squared, F-statistic.}
\entry{predict(model, newdata)}{Predicted \texttt{y} for the rows of \texttt{newdata}.}
\entry{summary(model)\$r.squared}{Pull out R-squared directly. Look for \texttt{Multiple R-squared} near the bottom of \texttt{summary(model)}.}

\section*{Logistic regression}
\entry{glm(y \textasciitilde{} x1 + x2, family = binomial, data = df)}{Fit a logistic regression.}
\entry{predict(model, newdata, type = "response")}{Predicted probabilities for the rows of \texttt{newdata}. Without \texttt{type = "response"} you get the linear predictor; pass it through \texttt{plogis()} to convert it to a probability.}

\section*{Train / validation / test split}

The index-sampling pattern:

\begin{verbatim}
n     <- nrow(df)
idx   <- sample(1:n, size = floor(0.7 * n))
train <- df[idx, ]
val   <- df[-idx, ]
\end{verbatim}

For a three-way split, sample a second index from the leftover rows.

\section*{Diagnostic plots}
\entry{plot(model)}{Cycles through residuals-vs-fitted, normal Q--Q, scale--location, and residuals-vs-leverage.}

\section*{Reading a regression table}

In \texttt{summary(model)}, each row of the \texttt{Coefficients} block has columns \texttt{Estimate}, \texttt{Std. Error}, \texttt{t value} (or \texttt{z value} for \texttt{glm}), and \texttt{Pr(>|t|)} (or \texttt{Pr(>|z|)}). An approximate 95\% confidence interval for a coefficient is \texttt{Estimate +/- 1.96 * Std. Error}. The line beginning \texttt{Multiple R-squared} near the bottom is the R-squared.

\end{multicols}
\end{document}
